\chapter{Introducción}
\label{cap:introduccion}

El mundo de la informática en general atraviesa una transición importante. Con la
explosión de la Inteligencia Artificial, procesar datos en tiempo real ha dejado de ser
un lujo para convertirse en una necesidad.
Si bien el enfoque tradicional pasaba por delegar el cómputo de las redes neuronales en
grandes centros de datos en la nube, hoy la tendencia apunta hacia el concepto de
Edge-AI: llevar la inteligencia directamente al extremo, al dispositivo local.
¿Qué logramos con esto? Reducir de forma drástica la latencia, blindar la privacidad del
usuario al eliminar el envío de información por la red y aliviar la carga de las redes de
comunicación \citep{zhou2019edgeintelligencepavingmile}.

Aquí es donde entran en juego las FPGA como pieza clave.
A diferencia de las CPU y GPU tradicionales que procesan la información de forma
secuencial o en paralelo pero de forma rígida, las FPGA permiten diseñar arquitecturas
de hardware como si fueran un lienzo en blanco, que optimizan el flujo de datos y el
paralelismo. Esto es importante para superar barreras físicas
actuales como el <<Power Wall>>, donde seguir subiendo la frecuencia de
los procesadores convencionales ha dejado de ser viable a nivel energético
\citep{10.1145/3282307}.

Por todo ello, este Trabajo de Fin de Grado se centra en explorar a fondo el co-diseño
hardware-software.
Para materializarlo utilizaremos la plataforma Digilent Zybo \citep{digilentZyboOverview},
una placa que combina la flexibilidad de un procesador ARM Cortex-A9 para tareas de
control con la potencia de una FPGA Xilinx/AMD de la familia 7-series para tareas de
cálculo intensivo en un mismo chip. A nivel académico es un entorno ideal para emular
sistemas industriales, dándole al diseñador el poder de controlar cada ciclo de reloj y
cada transferencia de datos a través de buses de alto rendimiento como AXI4.

\section{Motivación}
La demanda de capacidad de cómputo, impulsada principalmente por algoritmos de
Inteligencia Artificial y Deep Learning, crece a un ritmo que los procesadores de
propósito general ya no pueden satisfacer eficientemente bajo restricciones estrictas de
consumo. Trasladar la inferencia desde la nube hasta los equipos locales no es opcional,
es el camino a seguir.

El verdadero reto técnico de este trabajo, y lo que lo hace tan desafiante, es intentar
exprimir al máximo una FPGA con recursos modestos (apenas 80 bloques DSP) para hacer
correr modelos de Machine Learning. Todo esto apoyándonos en herramientas modernas de
Síntesis de Alto Nivel (HLS) para desmitificar y agilizar la histórica complejidad que
siempre ha rodeado al diseño puramente hardware.

\section{Objetivos}
El objetivo general de este trabajo es claro: diseñar, implementar y evaluar un sistema
empotrado heterogéneo basado en el SoC Zynq Z-7010 capaz de acelerar la inferencia de un
modelo ligero de Inteligencia Artificial.
Derivado a esto, se quiere cuantificar las mejoras de rendimiento en hardware respecto a
la ejecución del mismo flujo de inferencia sobre el microprocesador ARM.
Para alcanzar esta meta, se definen los siguientes objetivos específicos:

\begin{enumerate}
  \item \textbf{Puesta en marcha de la plataforma:} configurar desde cero el entorno de
    desarrollo Vivado/Vitis para la placa Digilent Zybo, validando el funcionamiento
    independiente del microprocesador ARM Cortex-A9 y de la FPGA.
  \item \textbf{Diseño de la infraestructura de comunicación:} implementar canales de
    comunicación de alto ancho de banda entre el Sistema de Procesamiento y la Lógica
    Programable, utilizando el protocolo AMBA AXI y controladores de Acceso Directo a Memoria (DMA).
  \item \textbf{Aceleración de un cálculo base:} escribir en \texttt{C++} y sintetizar a hardware
    (HLS) un IP core básico de álgebra lineal (multiplicación de matrices) para validar
    el flujo de datos entre procesador y FPGA.
  \item \textbf{Implementación de IA en recursos limitados:} adaptar un modelo de red
    neuronal (Red Neuronal Convolucional CIFAR-10) al conjunto reducido de recursos de la
    Zynq Z-7010, generando su hardware equivalente y desarrollando el software de control
    que le suministre los datos de entrada.
  \item \textbf{Evaluación de rendimiento (Benchmarking):} realizar una comparativa de
    tiempos de ejecución (latencia) lanzando la inferencia del modelo neuronal únicamente
    en la CPU frente a la ejecución asistida por la FPGA.
\end{enumerate}

\section{Plan de trabajo}
Para organizar el esfuerzo y cumplir los hitos, el proyecto se divide en estas fases incrementales:

\textbf{Fase 1.} Investigación y familiarización: revisión del estado del arte en sistemas Zynq,
buses AXI y Síntesis de alto nivel. Instalación de herramientas (Vivado, Vitis) y
ejecución de <<Hola Mundo>> en la placa.

\textbf{Fase 2.} Comunicación PS-PL: configuración de registros de control mediante AXI-Lite y
configuración del subsistema de memoria con AXI DMA. Desarrollo de drivers en C
(baremetal) para mover datos desde la RAM a la FPGA.

\textbf{Fase 3.} Desarrollo del acelerador base en HLS: programación en
\texttt{C}/\texttt{C++} de una
multiplicación de matrices. Uso de pragmas HLS (PIPELINE, UNROLL) para optimizar el
hardware. Integración del bloque IP generado en Vivado con el DMA.

\textbf{Fase 4.} Diseño e implementación del modelo de IA: entrenamiento y/o exportación de un
modelo de Red Neuronal Convolucional para clasificación de imágenes CIFAR-10. Generación
del acelerador hardware de la red neuronal y conexión en el diseño de Vivado.

\textbf{Fase 5.} Pruebas y benchmarking: ejecución de la inferencia en FPGA y comparación con
el core generado por hls4ml ejecutado en el ARM. Toma de métricas de rendimiento y verificación de
resultados.

\textbf{Fase 6.} Redacción de la memoria: documentación del proceso, análisis de los resultados
obtenidos y redacción final del documento del TFG.
